\documentclass{report}
\usepackage{amsmath,amssymb}
\setlength{\parindent}{0mm}
\setlength{\parskip}{1em}
\begin{document}
\begin{center}
\rule{6in}{1pt} \
{\large Geoffrey D. Sanders \\
{\bf Hub Snub: removing vertices with high degree from coarse-grid correction}}

7000 East Avenue \\ L-561 \\ Livermore \\ CA 94551
\\
{\tt sanders29@llnl.gov}\\
Alyson Fox\\
Thomas Manteuffel\\
Tobias Jones\end{center}

Network scientists often employ numerical solutions to linear systems as
subroutines of data mining algorithms. Due to the ill-conditioned nature
of the systems, obtaining solutions with standard iterative methods is
often prohibitively costly; current research aims to automatically
construct preconditioners that are optimally efficient and scalable for a
very broad class of graph associated matrices and network topologies.
Frequently, networks encountered have skewed degree distribution and
small-world topology and we present new strategies to aggressively
coarsen the underlying systems into significantly smaller,
better-conditioned systems, from which smooth error is approximated with
high accuracy. Aggressive coarsening strategies are often applicable as a
preprocessing step for almost any other preconditioner. We demonstrate
their efficacy in conjunction with adaptive algebraic multigrid.

In this talk, we focus on analyzing coarse grid selection strategies that
eliminate high-degree vertices. We observe that there is often relatively
little energy in smooth modes at central, high-degree vertices and
propose completely ignoring these hubs in coarse-grid representation.
Clever smoother design further improves this property. We develop an
adaptive approach to identify which hubs are best candidates for removal
from coarse spaces in the setup of the multilevel hierarchy. We discuss
the trade-off between convergence factors and complexity involved in this
process.


\end{document}
